%! The Language of Variation — Unit 1, Lesson 1.1 — Lesson Plan
% GRADUAL-RELEASE plan (warm-up / guided notes and practice / debrief / close and assign).
% THERE IS NO GROUP ACTIVITY IN THIS COURSE — the guided notes carry the entire release,
% I do -> we do -> you do, and the debrief closes them. Teacher-facing. Every box is
% `skillbox` (breakable) with a \boxguard ahead of it; this course has no `fixedskillbox`,
% no tier boxes, and no exit ticket.
\documentclass[10pt]{article}
\usepackage{apstats-article}
\usepackage{apstats-boxes}

\newcommand{\UnitNumberName}{Unit 1: Exploring One-Variable Data \quad}
\newcommand{\LessonNumberName}{Lesson 1.1: The Language of Variation}

\begin{document}
\begin{center}
    {\Large\bfseries \CourseName \\
    {\small\bfseries \UnitNumberName \LessonNumberName}}
\end{center}

\begin{tcolorbox}[colback=sky, colframe=navy, boxrule=0.9pt, arc=2mm,
    left=2mm, right=2mm, top=1.5mm, bottom=1.5mm]
    \textbf{Primary Objective:} Students will identify the variables recorded in a data set and
    classify each as categorical or quantitative --- judging by whether the values record a
    measured or counted amount rather than by whether they are written in digits.
    \\[2pt]
    \textbf{CED alignment:} Topic 1.2 \quad$\cdot$\quad Big Idea: VAR (Variation \& Distribution)
    \quad$\cdot$\quad Skill 2.A \quad$\cdot$\quad LO VAR-1.B, VAR-1.C \quad$\cdot$\quad
    EK VAR-1.B.1, VAR-1.C.1, VAR-1.C.2
    \\[2pt]
    \textbf{Lesson model:} \emph{Gradual release --- I do, we do, you do --- all of it inside
    the guided notes.} There is no group activity. The notes deliver both definitions and the
    amount test directly; the guided practice is worked together with students holding the pen;
    Independent Practice is where students carry it alone. The whole-class debrief is
    \emph{spoken} --- there is no transfer set and no reflection box on the page.
\end{tcolorbox}

\renewcommand{\arraystretch}{1.4}
\boxguard
\begin{skillbox}[Learning Targets \& Key Understandings]{goldbox}
\begin{tabularx}{\linewidth}{@{}X|X@{}}
\textbf{Learning targets (student language)} & \textbf{Key understandings (the ``why'')} \\[2pt]
\hline
\vspace{-6pt}
\begin{itemize}[leftmargin=1.1em, nosep, after=\vspace{2pt}]
  \item I can classify each variable in a data set as categorical or quantitative.
  \item I can run the amount test --- does the value record how much or how many? --- instead
        of asking whether it is written in digits.
  \item I can explain why averaging a ZIP code or a permit number is meaningless.
  \item I can explain what changes, and what is lost, when measured amounts are replaced by
        group labels.
\end{itemize}
&
\vspace{-6pt}
\begin{itemize}[leftmargin=1.1em, nosep, after=\vspace{2pt}]
  \item The classification decides everything downstream: which graph is legal, which summary
        is legal, and which inference procedure applies (EK VAR-1.C.1, VAR-1.C.2).
  \item \textbf{Target misconception:} that a column written in digits is quantitative.
        Permit numbers, ZIP codes, parking spots, jersey numbers, and PINs are all digits and
        all categorical --- their digits \emph{identify}, they do not \emph{measure}. The test
        is whether the value is an \emph{amount}.
  \item \textbf{Second misconception:} that producing a percent makes a variable quantitative.
        Counting is exactly how a categorical variable is summarized; the summary is numeric,
        the variable is not.
\end{itemize}
\end{tabularx}
\end{skillbox}

\boxguard
\begin{skillbox}[{Vocabulary, Concepts, \& Theorems --- taught directly in the guided notes}]{greenbox}
\begin{minipage}[t]{0.48\linewidth}
    \begin{itemize}
        \item \textbf{Variable} --- a characteristic that changes from one individual to
              another (EK VAR-1.B.1)
        \item \textbf{Categorical variable} --- takes on values that are category names or
              group labels (EK VAR-1.C.1)
        \item \textbf{Quantitative variable} --- takes on numerical values for a measured or
              counted quantity (EK VAR-1.C.2)
    \end{itemize}
\end{minipage}
\hfill
\begin{minipage}[t]{0.48\linewidth}
    \begin{itemize}
        \item \textbf{The amount test} --- ask whether the value records how much or how many,
              not whether it is written in digits
        \item \textbf{Grouping (binning)} --- replacing measured amounts with group labels;
              turns a quantitative variable into a categorical one, and the mean cannot be
              recovered
        \item \textbf{Carried forward from Lesson 1.0} --- individual, value, statistical
              question
    \end{itemize}
\end{minipage}
\end{skillbox}

% A tabularx never splits, so guard generously ahead of it.
\boxguard[30]
\begin{skillbox}[Lesson at a Glance --- \MeetingLength]{sky}
\renewcommand{\arraystretch}{1.25}
\begin{tabularx}{\linewidth}{@{}l c X X@{}}
\textbf{Phase} & \textbf{Min} & \textbf{Students} & \textbf{Teacher} \\
\hline
Warm-Up & 5 & An instrument rental sheet; three items, alone then with a neighbor & Debrief item~2 aloud; leave the ``which lists can you add up'' answer on the board \\
Guided Notes \& Practice & 38 & Fill the vocabulary box and notes 1--4; the Season Pass together; then Independent Practice alone & \textbf{I do} notes 1--4 (18) $\to$ \textbf{we do} the Guided Practice box (12) $\to$ \textbf{you do} Independent Practice (8) while you circulate \\
Debrief & 8 & Pens down, papers up --- answer aloud, nothing to fill in & Take the four digit-columns of notes~3 public and run the amount test on each aloud \\
Close \& Assign & 4 & Start the homework; hear the preview & Announce the paper homework and preview Lesson~1.2 \\
\end{tabularx}
\end{skillbox}

\boxguard[20]
\begin{skillbox}[Warm-Up --- Activate Prior Knowledge (5 min)]{sky}
\begin{minipage}[t]{0.48\linewidth}
    \textbf{The three items and what each seeds}
    \begin{itemize}
        \item \textbf{1.} Individuals, one variable, one of its values, on a rental sheet.
              \textbf{Seeds:} the Lesson~1.0 spiral --- and it is the vocabulary today's notes
              build on.
        \item \textbf{2.} ``Which of these lists could you add up and get something meaningful?''
              \textbf{Seeds:} the amount test, directly. Their own answer to ``what do the ones
              you circled have in common'' becomes notes section~2.
        \item \textbf{3.} The school ID number. \textbf{Seeds:} the digits trap, which is the
              crux of the whole lesson.
    \end{itemize}
\end{minipage}
\hfill
\begin{minipage}[t]{0.48\linewidth}
    \textbf{Running it}
    \begin{itemize}
        \item Item~2's answers are (A) prices and (C) ages. Expect some students to circle (B)
              or (D) --- jersey numbers and house numbers add up cleanly and mean nothing. That
              error \emph{is} the lesson; do not correct it in the warm-up, write it on the board.
        \item \textbf{Debrief item~2 aloud} and leave their phrasing up. You will point at it
              during notes section~2 and again in the closing debrief.
        \item Item~3 needs 30 seconds. ``Your ID number is not measuring anything about you ---
              it is naming you.'' Then stop; the notes do the rest.
    \end{itemize}
\end{minipage}
\end{skillbox}

\boxguard
\begin{skillbox}[Guided Notes \& Practice --- I do, we do, you do (38 min)]{sky}
\textbf{This is the whole lesson.} There is no group activity; the notes carry the entire
release, and Independent Practice is where students carry it alone.
\begin{multicols}{2}
\textbf{I do (about 18 min).} Present notes sections 1--4 in order; students fill the blanks as
each idea is named, and the vocabulary box is filled \emph{as you go}, never in advance.

\emph{Section 1} gives both definitions on one three-row table --- trace \emph{Trail} and
\emph{Time out} with a finger and say out loud that 215 really is 120 more than 95, while Summit
Trail is not ``more'' than Ridge Loop. \emph{Section 2} is the amount test, in a box, in one
sentence. \textbf{Section 3 is the lesson.} Work the \texttt{work} block on the board: both
averages compute, one is an amount and one is nobody's permit. Then release the five-column C/Q
strip --- expect the room to split on \emph{Parking spot}. Close section~3 on ``62\% of parties
hiked with a dog'': counting produces numbers, and the variable is still categorical.
\emph{Section 4} is grouping, and the sentence to land is \textbf{grouping runs one way only.}

\textbf{We do (about 12 min).} The Guided Practice box, \emph{The Season Pass}. Students hold the
pen; you hold the questions: \emph{``Pass number --- run the amount test on it out loud.''}
\quad \emph{``Miles traveled and locker number are both digits. Say the difference in one
sentence.''} \quad Part (b) is where grouping gets used rather than recited: insist on
\emph{what can no longer be computed}, not just ``it becomes categorical.''

\columnbreak
\textbf{You do (about 8 min).} \emph{Independent Practice}, three items, silent and alone.
\textbf{Item~1 is the formative check} --- the coffee-shop strip has one of each trap in it
(\emph{Drink size} is a label that looks ordered; \emph{Loyalty card \#} is digits that label).
Sort what you see into three piles: correct with a reason (ready); correct with no reason (ready,
needs the language); marked Q (the misconception survived). \textbf{If the third pile is more
than a third of the room, spend the debrief reteaching section~3} rather than harvesting reasons.

Circulate and demand the reason, not the letter:
\begin{itemize}[leftmargin=1.1em, nosep]
  \item \textbf{Item 1, the crux:} ``Average two loyalty card numbers. Now hand me the customer
        whose card that is.''
  \item \textbf{Item 1:} ``Large really is more than Small --- so how much more?''
  \item \textbf{Item 3:} ``What would a summary of \emph{Type of pet} even look like?''
\end{itemize}
While circulating, pick the papers to put up in the debrief --- one that got \emph{Loyalty
card \#} wrong, one that justified it correctly.
\end{multicols}
\end{skillbox}

\boxguard
\begin{skillbox}[Debrief --- whole class (8 min)]{sky}
Pens down, papers up. \textbf{This debrief is spoken} --- there is no box at the end of the
notes to fill in, so it lives entirely on the board and in the room.
\begin{multicols}{2}
\begin{enumerate}[leftmargin=1.3em, nosep, itemsep=3pt]
  \item \textbf{Put the four digit-columns from notes section~3 on the board and run the amount
        test on each aloud}, in this order: \emph{Water bottles} (easy, Q), \emph{Permit \#}
        (easy, C), \emph{Home ZIP} (C, and say ``80211 is not more than 80304''),
        \emph{Parking spot} (the hard one --- spot 17 is not more parking than spot 04).
  \item Put up the two Independent Practice papers you picked and name what the correct one did:
        \textbf{categorical variables are summarized by counting, not averaging.}
  \item Settle the percent trap in one line, off notes~3: the summary is numeric, the variable
        is not.
  \item Take grouping last. ``Long'' cannot be turned back into 215 minutes.
\end{enumerate}
\columnbreak
\textbf{Two things to demand aloud.} First, \textbf{make someone state the amount test without
the word ``digits''} --- that constraint is the point, because ``it has digits'' is the rule
students are carrying in. Second, \textbf{make them generate a pair of examples that were not on
any page today}, one Q and one C, with the reason attached; that is where transfer shows up now
that there is no transfer set on paper. Hold the room to both rather than letting students race
ahead to the homework.

\textbf{Connect back to Lesson~1.0 if time allows:} \emph{``Is `What kind of pet do students own?'
a statistical question?''} Categorical answer, and still statistical --- the same point the gym
ZIP code made yesterday. Students who classify ZIP as quantitative today are usually the ones who
called that question non-statistical yesterday.

\textbf{If the debrief runs short}, cut the Lesson~1.0 callback --- Debrief item~1 is the only
one that cannot be cut.
\end{multicols}
\end{skillbox}

\boxguard
\begin{skillbox}[AP Practice --- assigned, not scored]{goldbox}
Four AP-style multiple-choice items (five options each, in context) on a trailhead register ---
a data set students have not seen, and the only place it appears in the packet ---
then one multi-part free-response set on a city library's cardholder records: classify four
variables, explain why an average card number is not a summary, explain what grouping costs, and
--- part (d), the spiral --- write a statistical question that uses a \textbf{categorical}
variable.

\textbf{Not scored --- the cover's score column reads \textbf{NA} for this page.} Go over it at
the start of Lesson~1.2. MC item~2 is the ZIP trap, item~3 the percent trap, item~4 grouping.
Part (d) is the one worth reading closely: it is the Lesson~1.0 spiral, and a student who cannot
write a categorical statistical question has not separated \emph{varies} from \emph{numeric}.
\end{skillbox}

\boxguard
\begin{skillbox}[Homework --- scored, due next class]{goldbox}
Five items in fresh contexts: a coffee-shop order log (the C/Q strip), a pet survey that pairs
\emph{Number of pets} with \emph{Type of pet}, a hospital chart where a second column is
\emph{derived} from the first by grouping, the digits-that-label prompt, and one AP-style
multiple-choice item with a required justification. The packet closes with the Lesson~1.2
preview.

\textbf{Start it in the last four minutes of the period.} \textbf{Scoring:} 10 points --- 2 per
item. \textbf{Item~3 is the one to read closely}: the hospital chart is the only place students
meet a categorical column that was \emph{manufactured} from a quantitative one, and explaining
where the second column came from is the grouping idea in its transferable form.

\textbf{Paper homework is the default for this course} --- assign this page unless you have
decided otherwise. \textbf{DeltaMath override (occasional):} on the days you assign DeltaMath
instead, it replaces this page, you strike through the score cell on the cover, and this page
stays in the packet as extra practice. Never assign both.
\end{skillbox}

\boxguard
\begin{skillbox}[Watch For (while circulating)]{redbox}
\begin{itemize}[leftmargin=1.2em, nosep]
  \item \textbf{``It has digits, so it is quantitative''} (notes~3, IP~1 and~2, HW~1 and~4).
        The single misconception this lesson exists to break. Probe: \emph{``Average two of
        them. Now hand me the individual whose value that is.''}
  \item \textbf{Ordered labels read as amounts} (IP~1 and HW~1, \emph{Drink size} S/M/L). Large really is
        more than Small --- but there is no number of anything. Probe: \emph{``How much more?''}
  \item \textbf{Percent taken as proof of quantitative} (notes~3, IP~3, AP MC~3). Probe:
        \emph{``What are the values in that column --- Yes and No, or 62?''}
  \item \textbf{Grouping described but not costed} (notes~4, GP~(b), HW~3). ``It becomes
        categorical'' is half an answer. Probe: \emph{``Name one thing you can no longer
        compute.''}
  \item \textbf{Counted quantities called categorical} (\emph{Group size}, \emph{Number of
        shots}). A count is an amount. Probe: \emph{``Is six more than two?''}
\end{itemize}
Cold-call prompts: \emph{``Give me a column of digits that is categorical, and say what the
digits do.''} \quad \emph{``Name a variable you could group, and what you would lose.''}
\end{skillbox}

\boxguard
\begin{skillbox}[Close \& Assign (4 min)]{goldbox}
Hand the homework launch first --- five items on paper, scored, due next class --- and point out
that the AP Practice page is theirs to work and bring to review. Then close on the decision this lesson bought: \emph{``You now know which of your
columns can be averaged and which can only be counted. That single call decides every graph we
are allowed to draw for the rest of the year --- starting tomorrow.''} \textbf{Preview:}
Lesson~1.2, \emph{Representing a Categorical Variable with Graphs} --- frequency and relative
frequency tables, and how the same counts can be made to tell two very different stories.
\end{skillbox}

% ── Teacher notes — one per component, in packet order (four: there is no activity) ──
\begin{teachernote}[Warm-Up]
Item~2 is the whole warm-up. The intended answers are (A) prices and (C) ages; jersey numbers and
house numbers are the distractors, and \textbf{students who circle them are showing you the
misconception the lesson is built to break} --- so do not correct it here. Write both the correct
and incorrect circles on the board without judgment and say ``hold that.'' You will settle it in
notes section~3, with the permit-number computation, where the evidence does the work instead of
your authority. Item~1 is a 60-second spiral to Lesson 1.0; item~3 takes 30 seconds and needs no
debrief beyond one sentence.
\end{teachernote}

\begin{teachernote}[Guided Notes \& Practice]
\textbf{This one document is the lesson --- pace it 18 / 12 / 8, then 8 for the debrief.}
Sections 1 and 2 are quick: definitions and one boxed sentence. \textbf{Spend the time in
section 3}, which carries the misconception. Work the \texttt{work} block live: compute both
averages, then ask for the party who holds permit 4179.5. The silence is the argument.

The five-column C/Q strip in section~3 is where you find out who is with you. \emph{Parking spot}
is the one that splits a room, because two-digit spot numbers look like small counts. Take a hand
count before revealing.

\textbf{Guided Practice is where the pen changes hands.} \emph{Locker \#} and \emph{Miles
traveled} are deliberately adjacent --- both digits, opposite answers. Part (b) must end with a
named loss (``the mean age''), not just a classification.

\textbf{Independent Practice is the formative check, and it is the last thing on the page ---
there is no transfer set and no reflection box after it.} Item~1's \emph{Loyalty card \#} tells
you whether section~3 landed. Sort as you circulate: correct with a reason (ready); correct with
no reason (ready, needs the language); marked Q (the misconception survived). If the third pile
is more than a third of the room, \textbf{spend the debrief reteaching section~3} instead of
harvesting reasons --- and say so plainly, because the AP Practice page behind the notes is where
the transfer now happens, on their own time.

The ground rule --- \emph{no letter goes down until you can say why} --- is worth enforcing for
the first two minutes even if it slows them down. The reasons are what the debrief harvests.
Students who finish early get the extension question: \emph{``give me a column of digits that
would behave the same way as Loyalty card \#.''}

\textbf{Debrief (8 min), spoken.} The four-column amount test from notes~3, in the order given
--- easy, easy, hard, hardest --- is the sequence that makes \emph{Parking spot} land. Do not
start with the hard one. If time is short, cut the Lesson~1.0 callback; making a student state
the test \emph{without the word ``digits''} is the one move that cannot be cut.
\end{teachernote}

\begin{teachernote}[AP Practice]
\textbf{This page is not scored} --- the graded practice for this lesson is the homework that
follows it. Work it as AP-format reps and go over it at the start of Lesson~1.2. \textbf{This page now
carries the transfer:} the trailhead register appears nowhere else in the packet, so the four
multiple-choice items are the first time students meet this data set --- treat what comes back
in Lesson~1.2 as your real evidence that the amount test transferred.

MC item~2 is the ZIP trap in AP wording, item~3 the percent trap, item~4 grouping. On the free
response, part~(b) requires the student to say what a card number \emph{does} instead of measure;
``it is categorical'' alone is not the answer. \textbf{Part (d) is the spiral} and the one to read
first --- it reaches back to Lesson~1.0 and makes students run the statistical-question test on a
categorical variable, which is exactly where the ``statistical $=$ numeric'' misconception hides.
\end{teachernote}

\begin{teachernote}[Homework]
\textbf{Scored, 10 points (2 per item), due next class.} Launch it in the last four minutes so
students hit item~1 while you can still answer.

\textbf{Item~3 is the one to read closely.} The hospital chart is the only place in the lesson
where students meet a categorical column that was \emph{manufactured} from a quantitative one:
exact body temperature in $^\circ$C, then a second column reading \emph{Normal} or \emph{Fever}.
A student who says only ``the first is Q and the second is C'' has classified but not understood;
the credit is in explaining that the second column was made \emph{from} the first by grouping,
and that the original temperatures cannot be recovered from it. Item~2 (\emph{Number of pets}
vs.\ \emph{Type of pet}) is the same distinction in its cleanest form and is worth using as the
warm-up spiral for Lesson~1.2.

\textbf{Paper is the default.} This page is what gets assigned unless you decide otherwise on a
given day. On the occasions you assign DeltaMath in its place, strike the score cell on the cover
and tell students this page is now optional practice. Do not assign both.
\end{teachernote}
\end{document}
